% !PRINTED 380
\weberpage{380}
\S\,104. \gesperrt{Characters and Genera of Quadratic Forms.}

By a fundamental-discriminant divisor $\delta$ of a discriminant $D$ we shall understand
the following (\S\,84):

1. \gesperrt{$\delta$ is a fundamental discriminant dividing $D$.}

2. \gesperrt{The quotient $D:\delta = D_1$ is itself again a discriminant.}

An odd fundamental-discriminant divisor can contain no prime factors other than
such as divide $D$, and none of them more than once; on the other
hand, any product of distinct odd primes dividing $D$, furnished
with a sign such that $\delta \equiv 1 \pmod{4}$, is always a fundamental-discriminant divisor.

Besides these, there occur among the fundamental-discriminant divisors:

\weq{1}{\begin{aligned}
&-4, \text{ if } D \equiv 0,\ -4 \pmod{16}, \\
&+8, \text{ \glqq } D \equiv 0,\ 8 \pmod{32}, \\
&-8, \text{ \glqq } D \equiv 0,\ -8 \pmod{32},
\end{aligned}}

If, then, we denote the odd fundamental-discriminant divisors by $\delta_1$, all the
fundamental-discriminant divisors $\delta$ come out as follows:
% [sic] Weber reuses the symbol delta_1 here for the odd stem-divisors, distinct from the delta_1 of the abstract-group discussion further below on this page.

\[
\begin{aligned}
1)\ & \delta = \delta_1, & D &\equiv\ \ \ 1 \pmod{4}, \\
& & &\equiv\ \ \ 4 \pmod{16}, \\
2)\ & \delta = \delta_1,\ -4\delta_1, & D &\equiv -4 \pmod{16}, \\
& & &\equiv\ 16 \pmod{32}, \\
3)\ & \delta = \delta_1,\ +8\delta_1, & D &\equiv\ \ \ 8 \pmod{32}, \\
4)\ & \delta = \delta_1,\ -8\delta_1, & D &\equiv -8 \pmod{32}, \\
5)\ & \delta = \delta_1,\ -4\delta_1, & D &\equiv\ \ \ 0 \pmod{32}. \\
& \ \ \ \ 8\delta_1,\ -8\delta_1, & &
\end{aligned}
\]

The number of fundamental-discriminant divisors (1 being counted as a fundamental-discriminant divisor)
is accordingly always a power of 2.

3. \gesperrt{If we set it equal to $2^\lambda$, then in the cases 1), 2),
3), 4) $\lambda$ is equal to the number of distinct primes
dividing $D$, and in the case 5) greater by one.}

By applying a symbolic multiplication, the fundamental-discriminant divisors of $D$
can be made into an \gesperrt{Abelian group}.

For if $\delta_1$, $\delta_2$ are two fundamental-discriminant divisors of $D$, then
the product $\delta_1\delta_2$ is likewise a discriminant, but not always
a fundamental discriminant. We denote by $\delta$ the \gesperrt{fundamental discriminant associated with}
of $\delta_1\delta_2$ and set symbolically

\weq{2}{\delta = \delta_1\delta_2.}
